% ============================================================
%  APUNTES ANGULAR — Survey Maker
%  Documento maestro — todos los bloques aqui
%  Compilar con: LuaLaTeX
%  Comando: lualatex ApuntesAngular.tex
% ============================================================
\documentclass[12pt, a4paper]{article}

% --- Fuentes ------------------------------------------------
\usepackage{fontspec}
\setmainfont{Latin Modern Roman}
\setmonofont[Scale=0.88]{Courier New}

% --- Emojis -------------------------------------------------
\usepackage{emoji}
\setemojifont{Segoe UI Emoji}

% --- Idioma -------------------------------------------------
\usepackage[spanish]{babel}

% --- Márgenes -----------------------------------------------
\usepackage[top=2.5cm, bottom=2.5cm, left=2.8cm, right=2.8cm]{geometry}

% --- Colores ------------------------------------------------
\usepackage[dvipsnames, table]{xcolor}
\definecolor{azulSignal}{RGB}{37, 99, 235}
\definecolor{verdeOk}{RGB}{22, 163, 74}
\definecolor{rojoError}{RGB}{220, 38, 38}
\definecolor{grisClaro}{RGB}{243, 244, 246}
\definecolor{grisCodigo}{RGB}{31, 41, 55}
\definecolor{amarillo}{RGB}{251, 191, 36}
\definecolor{morado}{RGB}{124, 58, 237}
\definecolor{naranja}{RGB}{234, 88, 12}
\definecolor{azulParte}{RGB}{15, 52, 96}

% --- Código fuente ------------------------------------------
\usepackage{listings}
\lstdefinestyle{ts}{
  language        = [Sharp]C,
  basicstyle      = \ttfamily\small,
  keywordstyle    = \color{azulSignal}\bfseries,
  commentstyle    = \color{verdeOk}\itshape,
  stringstyle     = \color{rojoError},
  backgroundcolor = \color{grisClaro},
  frame           = single,
  framerule       = 0.5pt,
  rulecolor       = \color{gray!40},
  numbers         = left,
  numberstyle     = \tiny\color{gray},
  numbersep       = 8pt,
  tabsize         = 2,
  breaklines      = true,
  showstringspaces = false,
  xleftmargin     = 1.5em,
  morekeywords    = {signal, computed, effect, inject, readonly,
                     const, let, async, of, map, return, true, false,
                     null, undefined, string, number, boolean, void,
                     update, untracked, input, output, model,
                     if, else, for, switch, case, default,
                     defer, placeholder, loading, error, empty,
                     track, on, when, ngOnInit, ngOnDestroy},
}
\lstset{style=ts}

% --- Cajas decorativas --------------------------------------
\usepackage[most]{tcolorbox}
\tcbset{
  analogia/.style={
    colback=blue!5, colframe=azulSignal,
    fonttitle=\bfseries\large, arc=6pt,
    left=8pt, right=8pt, top=6pt, bottom=6pt,
  },
  importante/.style={
    colback=yellow!10, colframe=amarillo!80!orange,
    fonttitle=\bfseries\large, arc=6pt,
    left=8pt, right=8pt,
  },
  regla/.style={
    colback=green!5, colframe=verdeOk,
    fonttitle=\bfseries\large, arc=6pt,
    left=8pt, right=8pt,
  },
  peligro/.style={
    colback=red!5, colframe=rojoError,
    fonttitle=\bfseries\large, arc=6pt,
    left=8pt, right=8pt,
  },
}

% --- Links --------------------------------------------------
\usepackage{hyperref}
\hypersetup{colorlinks=true, linkcolor=azulSignal, urlcolor=morado}

% --- Tablas -------------------------------------------------
\usepackage{booktabs}
\usepackage{array}

% --- Encabezado y pie ---------------------------------------
\usepackage{fancyhdr}
\pagestyle{fancy}
\fancyhf{}
\lhead{\textcolor{gray}{\small \emoji{open-book} ApuntesAngular}}
\rhead{\textcolor{gray}{\small Survey Maker}}
\cfoot{\textcolor{gray}{\thepage}}

% --- Separadores de Parte -----------------------------------
\usepackage{titlesec}
\titleformat{\part}[display]
  {\centering\normalfont\Huge\bfseries\color{azulParte}}
  {\centering\normalfont\large\color{gray} Parte \thepart}
  {12pt}{}

% --- Párrafos -----------------------------------------------
\setlength{\parskip}{6pt}
\setlength{\parindent}{0pt}

% ============================================================
\begin{document}

% --- Portada ------------------------------------------------
\begin{titlepage}
  \centering
  \vspace*{2cm}
  {\fontsize{56}{56}\selectfont \emoji{graduation-cap}}\\[0.5cm]
  {\Huge\bfseries\color{azulParte} Apuntes de Angular}\\[0.4cm]
  {\large\color{grisCodigo} Angular 17--19 · Signal-First · Survey Maker}\\[0.8cm]
  \textcolor{gray!50}{\rule{12cm}{1pt}}\\[0.8cm]

  \begin{minipage}{0.8\textwidth}
    \centering
    \begin{tcolorbox}[colback=grisClaro, colframe=gray!40, arc=6pt]
      \centering
      \textbf{Contenido}\\[0.4cm]
      \begin{tabular}{cl}
        \emoji{zap}           & Signals (\texttt{signal}, \texttt{computed}, \texttt{effect}) \\[4pt]
        \emoji{puzzle-piece}  & Bloque 1 — Componentes (\texttt{@Component}, \texttt{input}, \texttt{output}) \\[4pt]
        \emoji{artist-palette}& Bloque 2 — Templates (\texttt{@if}, \texttt{@for}, \texttt{@defer}) \\[4pt]
        \emoji{wrench}        & (próximos bloques aquí) \\
      \end{tabular}
    \end{tcolorbox}
  \end{minipage}

  \vspace{1.2cm}
  {\normalsize \emoji{student} \textbf{Aldo Zetina} — UNAM Ciencias de la Computación}\\[0.3cm]
  {\normalsize \texttt{\textcolor{morado}{github.com/AldoZM/Survey-Maker}}}\\[0.3cm]
  {\normalsize \emoji{calendar} Mayo 2026}\\[1.5cm]

  \begin{tcolorbox}[importante, title={\emoji{wrench} Compilar localmente}]
    \texttt{lualatex ApuntesAngular.tex}\\
    (requiere MiKTeX o TeX Live con LuaLaTeX)
  \end{tcolorbox}
\end{titlepage}

\tableofcontents
\newpage

% ============================================================
% ============================================================
\part{\emoji{zap} Signals — El Sistema Reactivo}
% ============================================================
% ============================================================

\section{\emoji{brain} La Idea Central — Sin Código}

Antes de ver código, entiende el concepto con una historia.

\begin{tcolorbox}[analogia, title={\emoji{magic-wand} La Pizarra Mágica}]
Imagina que tienes una \textbf{pizarra mágica} en tu cuarto.
Escribes un número en ella: \textit{``0 respuestas contestadas''}.

Afuera hay tres amigos mirando por la ventana:
\begin{itemize}
  \item \emoji{bar-chart} \textbf{Barra de progreso} — actualiza su dibujo cuando cambia el número.
  \item \emoji{white-check-mark} \textbf{Botón de Enviar} — se activa cuando el número llega al máximo.
  \item \emoji{floppy-disk} \textbf{localStorage} — guarda una foto del número automáticamente.
\end{itemize}
Cada vez que \textbf{borras el número y escribes uno nuevo}, los tres amigos lo ven al instante.
Tú no tienes que avisarles — ellos solos reaccionan.\\[0.3cm]
\textbf{Eso es un Signal.} La pizarra mágica.
\end{tcolorbox}

\section{\emoji{pencil} Tipo 1 — \texttt{signal()} — La Pizarra Mágica}

\subsection{Qué es}
Un \texttt{signal()} crea una \textbf{variable reactiva}. Cuando su valor cambia,
Angular actualiza todo lo que depende de él automáticamente.

\subsection{\emoji{file-folder} En el proyecto}
Archivo: \textcolor{morado}{\texttt{core/services/survey-state.service.ts}}, línea 28.

\begin{lstlisting}
// signal() crea la "pizarra magica".
// Todo empieza vacio {} — nadie ha contestado nada todavia.
// Record<string, FieldAnswer> = diccionario donde:
//   la LLAVE es el ID del campo  (ej: "nombre", "calificacion")
//   el VALOR es la respuesta     (ej: { value: "Juan", isValid: true })
readonly answers = signal<Record<string, FieldAnswer>>({});
//                                                      ^
//                                    {} = valor inicial, pizarra vacia
\end{lstlisting}

\subsection{\emoji{books} Cómo leer y escribir}

\begin{lstlisting}
// LEER — siempre con parentesis ()
const todasLasRespuestas = this.answers();
//                                      ^
//                    () = "dame lo que hay en la pizarra ahora"

// ESCRIBIR con .set() — reemplaza TODO el valor
this.answers.set({ nombre: { value: "Juan", isValid: true } });

// ACTUALIZAR con .update() — modifica sin perder lo anterior
this.answers.update(prev => ({
//                    ^
//             prev = lo que habia antes en la pizarra
  ...prev,           // copia todo lo anterior (no lo borres)
  [fieldId]: {       // solo cambia ESTE campo especifico
    value: "nuevo valor",
    isValid: true,
    isDirty: true,
  }
}));
\end{lstlisting}

\begin{tcolorbox}[regla, title={\emoji{pushpin} Regla de Oro}]
  \texttt{this.answers} \textrightarrow\ objeto Signal (no el valor).\\
  \texttt{this.answers()} \textrightarrow\ el valor real (el diccionario).\\[0.2cm]
  \textbf{Siempre leer con \texttt{()}.}
\end{tcolorbox}

\section{\emoji{abacus} Tipo 2 — \texttt{computed()} — El Amigo Matemático}

\subsection{Qué es}
Un \texttt{computed()} es una variable que \textbf{se calcula sola} a partir de
otros Signals. Cuando cualquier Signal que lee cambia, el \texttt{computed}
recalcula automáticamente. Tú nunca tienes que llamarlo.

\begin{tcolorbox}[analogia, title={\emoji{memo} Celda de Excel con Fórmula}]
  Imagina \texttt{C1 = A1 + B1} en Excel.\\
  Cuando cambias \texttt{A1} o \texttt{B1}, la celda \texttt{C1} se actualiza sola.\\[0.3cm]
  \emoji{zap} \texttt{progress} es \texttt{C1}. \quad \texttt{answers} es \texttt{A1}.
\end{tcolorbox}

\subsection{\emoji{bar-chart} Computed de Progreso}
Archivo: \textcolor{morado}{\texttt{core/services/survey-state.service.ts}}, línea 43.

\begin{lstlisting}
readonly progress = computed(() => {
  const s = this.schema();
  if (!s) return 0;

  // Aplanamos todas las secciones a un solo array de campos
  const allFields = s.sections.flatMap(sec => sec.fields);
  const requiredFields = allFields.filter(f => f.validation?.required);
  if (requiredFields.length === 0) return 100;

  // Un campo cuenta SOLO si: isValid=true Y isDirty=true
  const answeredRequired = requiredFields.filter(f => {
    const ans = this.answers()[f.id];
    //               ^
    //   LEE "answers" Signal -> crea la dependencia magica.
    //   Si "answers" cambia -> este computed recalcula solo.
    return ans?.isValid && ans.isDirty;
  });

  // (contestados / total) x 100
  return Math.round((answeredRequired.length / requiredFields.length) * 100);
});
\end{lstlisting}

\subsection{\emoji{white-check-mark} Computed de Validez}
Archivo: \textcolor{morado}{\texttt{core/services/survey-state.service.ts}}, línea 72.

\begin{lstlisting}
// ?puede el usuario enviar la encuesta?
readonly isValid = computed(() =>
  // .every() = "todos los campos cumplen la condicion?"
  // Si UN SOLO campo falla -> regresa false -> boton bloqueado
  this.schema()?.sections
    .flatMap(s => s.fields)
    .every(f => {
      if (!f.validation?.required) return true; // opcional -> valido
      const ans = this.answers()[f.id]; // dependencia magica
      return ans?.isValid ?? false;
    }) ?? false
);
\end{lstlisting}

\section{\emoji{shield} Tipo 3 — \texttt{effect()} — El Guardián Automático}

\subsection{Qué es}
Un \texttt{effect()} ejecuta código cada vez que un Signal cambia.
No retorna nada — solo realiza \textbf{efectos secundarios}: guardar datos, hacer log, llamar APIs.

\begin{tcolorbox}[analogia, title={\emoji{alarm-clock} La Alarma Automática}]
  La configuras una vez y te olvidas.\\
  Se activa sola cada vez que el Signal cambia.\\[0.3cm]
  En el proyecto: \textit{``cada vez que \texttt{answers} cambie, guárdalo en localStorage''.}
\end{tcolorbox}

Archivo: \textcolor{morado}{\texttt{core/services/survey-state.service.ts}}, línea 88.

\begin{lstlisting}
effect(() => {
  // untracked() = "lee este Signal pero NO lo vigiles"
  // Solo queremos reaccionar cuando ANSWERS cambia, no cuando schema cambia
  const s = untracked(() => this.schema());

  const answers = this.answers(); // ESTA linea crea la dependencia

  if (!s) return;
  try {
    // Autoguardado — como Google Docs guardando tu documento
    localStorage.setItem(`survey-state-${s.id}`, JSON.stringify(answers));
  } catch {
    // localStorage puede fallar en modo privado o si esta lleno
  }
}, { injector: this.injector });
\end{lstlisting}

\section{\emoji{eye} Tipo 4 — \texttt{computed()} para Lógica Condicional}

Decide si un campo se muestra según la respuesta de otro campo.
Archivo: \textcolor{morado}{\texttt{features/survey/field-renderer.component.ts}}, línea 156.

\begin{lstlisting}
// "el campo de comentarios solo aparece si elegiste 'Malo'"
readonly isVisible = computed(() => {
  const condition = this.field().condition;
  if (!condition) return true; // sin condicion -> siempre visible

  const answer = this.answers()[condition.dependsOn];
  const expected = condition.equals;

  // equals puede ser array: ["malo", "regular"] -> cualquiera vale
  if (Array.isArray(expected)) return expected.includes(answer as string);
  return answer === expected;
});
\end{lstlisting}

\begin{tcolorbox}[importante, title={\emoji{exclamation-mark} No es \texttt{display:none}}]
  Cuando \texttt{isVisible()} = \texttt{false}, Angular \textbf{elimina el componente del DOM}.\\
  No existe en memoria — cero costo de rendimiento.
\end{tcolorbox}

\section{\emoji{no-entry} Inmutabilidad — Por Qué Nunca Mutamos Directo}

\begin{lstlisting}
// MAL — misma referencia, Angular NO detecta el cambio
prev[fieldId] = nuevoValor;
this.answers.set(prev); // Angular ve el mismo objeto -> no re-renderiza

// BIEN — nueva referencia, Angular SÍ detecta el cambio
this.answers.update(prev => ({
  ...prev,               // copia todos los campos anteriores
  [fieldId]: nuevoValor  // sobreescribe solo el que cambio
}));
\end{lstlisting}

\begin{tcolorbox}[analogia, title={\emoji{house} Por Qué Importa la Referencia}]
  Angular compara \textbf{direcciones de objetos}, no su contenido.\\
  Con \texttt{\{...prev\}} creas un objeto \textit{nuevo} en otra dirección de memoria \textrightarrow\ Angular detecta el cambio y actualiza la UI.
\end{tcolorbox}

\section{\emoji{clipboard} Resumen — Signals}

\begin{center}
\renewcommand{\arraystretch}{1.6}
\begin{tabular}{>{\bfseries}p{3cm} p{4.5cm} p{5cm}}
\toprule
Tipo & Para qué & Ejemplo en el proyecto \\
\midrule
\texttt{signal()} & Guardar estado que cambia & \texttt{answers}, \texttt{schema}, \texttt{isSubmitted} \\
\texttt{computed()} & Calcular a partir de Signals & \texttt{progress}, \texttt{isValid}, \texttt{isVisible} \\
\texttt{effect()} & Efectos secundarios automáticos & Guardar en \texttt{localStorage} \\
\bottomrule
\end{tabular}
\end{center}

\newpage
% ============================================================
% ============================================================
\part{\emoji{puzzle-piece} Bloque 1 — El Componente a Fondo}
% ============================================================
% ============================================================

\section{\emoji{building-construction} ¿Qué es un Componente?}

Un componente es \textbf{la unidad mínima de UI en Angular}. Todo lo que ves
en pantalla es un componente — o varios componentes anidados.

\begin{tcolorbox}[analogia, title={\emoji{remote-control} El Control Remoto}]
Un componente tiene:
\begin{itemize}
  \item Una \textbf{carcasa} — el template HTML (lo que el usuario ve)
  \item \textbf{Botones que aceptan instrucciones} — los \texttt{input()}
  \item \textbf{Una pantalla que avisa hacia afuera} — los \texttt{output()}
  \item \textbf{Lógica interna} — la clase TypeScript
\end{itemize}
\end{tcolorbox}

\section{\emoji{microscope} Anatomía de \texttt{@Component}}

\begin{lstlisting}
@Component({
  // 1. Nombre del tag HTML que usas en templates
  //    Convencion: siempre "app-" + nombre-en-kebab-case
  selector: 'app-text-field',

  // 2. Sin NgModule (Angular moderno — el componente se autogestiona)
  standalone: true,

  // 3. Que otros componentes/directivas usa SU template
  //    Solo lo que usa ESTE componente, nada mas
  imports: [MatFormFieldModule, MatInputModule],

  // 4. El HTML del componente (inline o en archivo separado)
  template: `
    <mat-form-field>
      <input matInput [value]="value()" />
    </mat-form-field>
  `,

  // 5. CSS encapsulado — solo aplica a ESTE componente
  //    No choca con estilos de otros componentes
  styles: ['.full-width { width: 100%; }'],
})
export class TextFieldComponent { }
\end{lstlisting}

\section{\emoji{inbox-tray} \texttt{input()} — Recibir Datos del Padre}

\subsection{Viejo vs Nuevo}
\begin{lstlisting}
// VIEJO (Angular < 17) — decorador @Input()
@Input() value: string = '';
@Input({ required: true }) field!: QuestionField;

// NUEVO (Angular 17+) — funcion input()
readonly value = input<string>('');               // con default
readonly field = input.required<QuestionField>(); // obligatorio
\end{lstlisting}

\subsection{Por qué \texttt{readonly}}
\begin{lstlisting}
// readonly = el componente NO puede cambiar su propio input.
// Solo el PADRE puede cambiar el valor pasandole uno nuevo.
// Si intentas hacer this.field.set(...) -> error de TypeScript.
readonly field = input.required<QuestionField>();
\end{lstlisting}

\subsection{Los tres sabores}
\begin{lstlisting}
// 1. Opcional con default
readonly value = input<string>('');

// 2. Obligatorio — TypeScript error si el padre no lo pasa
readonly field = input.required<QuestionField>();

// 3. Con transform — transforma el valor antes de guardarlo
readonly disabled = input<boolean, string>(false, {
  transform: (v) => v === 'true' || v === true
});
\end{lstlisting}

\subsection{Cómo se ve desde el padre}
\begin{lstlisting}
<app-field-renderer
  [field]="field"                          <!-- [] = pasar dato al hijo -->
  [answers]="answersAsValues()"
  [value]="state.answers()[field.id]?.value"
  [errorMessage]="state.getError(field.id)"
  (valueChange)="state.setAnswer(field.id, $event)"  <!-- () = escuchar al hijo -->
/>
\end{lstlisting}

\section{\emoji{outbox-tray} \texttt{output()} — Avisar al Padre}

\subsection{Viejo vs Nuevo}
\begin{lstlisting}
// VIEJO — @Output() con EventEmitter
@Output() valueChange = new EventEmitter<string>();
this.valueChange.emit('nuevo valor');

// NUEVO — output()
readonly valueChange = output<string>();
this.valueChange.emit('nuevo valor');
\end{lstlisting}

\subsection{En \texttt{text-field.component.ts}}
\begin{lstlisting}
export class TextFieldComponent {
  readonly field = input.required<QuestionField>(); // recibe del padre
  readonly value = input<string>('');               // recibe del padre
  readonly valueChange = output<string>();           // avisa al padre

  onInput(event: Event): void {
    // Cada vez que el usuario escribe, avisamos al padre
    // $event en el template = lo que emitimos aqui
    this.valueChange.emit((event.target as HTMLInputElement).value);
  }
}
\end{lstlisting}

\begin{tcolorbox}[regla, title={\emoji{pushpin} Regla de los Corchetes}]
  \textbf{Datos bajan con \texttt{[ ]}} --- del padre al hijo.\\
  \textbf{Eventos suben con \texttt{( )}} --- del hijo al padre.\\[0.3cm]
  Juntos: \texttt{[( )]} = two-way binding (los dos en uno).
\end{tcolorbox}

\section{\emoji{syringe} \texttt{inject()} — Dependency Injection Moderna}

\subsection{Qué es}
\begin{tcolorbox}[analogia, title={\emoji{electric-plug} El Enchufe Central}]
  Angular tiene un ``enchufe central'' con todos los servicios disponibles.\\
  \texttt{inject()} es el cable que conecta tu componente a ese enchufe.\\
  Tú no fabricas el servicio — Angular te lo provee listo para usar.
\end{tcolorbox}

\subsection{Viejo vs Nuevo}
\begin{lstlisting}
// VIEJO — constructor injection
export class SurveyComponent {
  constructor(
    private state: SurveyStateService,
    private loader: SurveyLoaderService,
  ) { }
}

// NUEVO — inject() funcional
export class SurveyComponent {
  private readonly state  = inject(SurveyStateService);
  private readonly loader = inject(SurveyLoaderService);
}
\end{lstlisting}

\subsection{\texttt{providedIn: 'root'} — El Singleton}
\begin{lstlisting}
// Una SOLA instancia en toda la app
@Injectable({ providedIn: 'root' })
export class SurveyStateService { }
\end{lstlisting}

\begin{lstlisting}
// Todos obtienen EL MISMO objeto:
// SurveyComponent       inject(SurveyStateService) -> mismo objeto
// FieldRendererComponent inject(SurveyStateService) -> mismo objeto
// SurveyPageComponent   inject(SurveyStateService) -> mismo objeto

// Por eso cuando SurveyComponent guarda una respuesta,
// FieldRendererComponent la ve inmediatamente.
\end{lstlisting}

\section{\emoji{hourglass-done} Lifecycle Hooks — Los Momentos de Vida}

\begin{lstlisting}
// NACIMIENTO          VIDA                  MUERTE
// constructor()  ->   ngOnChanges() xN  ->  ngOnDestroy()
// ngOnInit() x1  ->   ngAfterViewInit() x1
\end{lstlisting}

\subsection{\texttt{ngOnInit()} — El más usado}
\begin{lstlisting}
export class SurveyPageComponent implements OnInit {
  readonly loader = inject(SurveyLoaderService);

  ngOnInit(): void {
    // Corre UNA SOLA VEZ cuando el componente aparece en el DOM.
    // Aqui van: llamadas HTTP iniciales, inicializacion de estado.
    // Los input() YA tienen sus valores en este punto.
    this.loader.loadSurvey('ejemplo-feedback');
  }
}
\end{lstlisting}

\subsection{\texttt{ngOnDestroy()} — Limpieza}
\begin{lstlisting}
export class MiComponente implements OnDestroy {
  private subscription = interval(1000).subscribe(() => console.log('tick'));

  ngOnDestroy(): void {
    // Corre cuando el componente se elimina del DOM.
    // Sin limpieza -> memory leak (la app se come la RAM).
    this.subscription.unsubscribe();
  }
}
\end{lstlisting}

\subsection{\texttt{effect()} como alternativa moderna a \texttt{ngOnChanges()}}
\begin{lstlisting}
// survey.component.ts — forma moderna con Signals
export class SurveyComponent {
  readonly schema = input.required<SurveySchema>();

  constructor() {
    // effect() detecta cuando schema() cambia — igual que ngOnChanges
    // pero con Signals: mas preciso, mas eficiente
    effect(() => {
      const s = this.schema();
      if (s) this.state.initSurvey(s);
    });
  }
}
\end{lstlisting}

\section{\emoji{link} \texttt{model()} — Two-Way Binding con Signals}

\begin{lstlisting}
// Sin model() — necesitas dos propiedades separadas
readonly value       = input<string>('');
readonly valueChange = output<string>();

// Con model() — una sola propiedad, two-way binding automatico
readonly value = model<string>('');

// En el template del padre:
// [(value)]="miVariable"
// Cuando el hijo hace value.set('x') -> miVariable se actualiza sola
\end{lstlisting}

\section{\emoji{clipboard} Resumen — Bloque 1}

\begin{center}
\renewcommand{\arraystretch}{1.6}
\begin{tabular}{>{\ttfamily\bfseries}p{4cm} p{8.5cm}}
\toprule
Concepto & Para qué \\
\midrule
@Component & Define un componente (selector, template, styles, imports) \\
input() & Recibir datos del padre (solo lectura) \\
input.required() & Input obligatorio — error si el padre no lo pasa \\
output() & Avisar al padre con eventos \\
model() & Two-way binding (input + output en uno) \\
inject() & Obtener servicios sin constructor \\
providedIn: 'root' & Singleton — una instancia para toda la app \\
ngOnInit() & Código al nacer (después de que inputs están listos) \\
ngOnDestroy() & Limpieza al morir (subscriptions, timers) \\
effect() & Reaccionar a cambios de inputs (forma moderna) \\
\bottomrule
\end{tabular}
\end{center}

\newpage
% ============================================================
% ============================================================
\part{\emoji{artist-palette} Bloque 2 — Templates Modernos (Angular 17+)}
% ============================================================
% ============================================================

\section{\emoji{brain} ¿Qué es un Template?}

El template es el HTML que define \textbf{lo que el usuario ve}.
Pero no es HTML normal --- Angular lo procesa y lo hace dinámico.

\begin{tcolorbox}[analogia, title={\emoji{cook} La Receta con Instrucciones}]
Piensa en el template como una \textbf{receta con instrucciones condicionales}:

\textit{``Si hay huevos, haz omelette. Si no hay, haz cereal.
Para cada ingrediente en la lista, agrégalo al plato.''}

Eso es exactamente lo que hacen \texttt{@if} y \texttt{@for}.
\end{tcolorbox}

\section{\emoji{speech-balloon} Interpolación \texttt{\{\{ \}\}}}

\begin{lstlisting}
<mat-card-title>{{ schema().title }}</mat-card-title>
<!--                ^  Signal — () obligatorios para leer -->

<span>{{ state.progress() }}% completado</span>
<!-- Si progress() = 60 renderiza: "60% completado" -->

<button>{{ schema().submitLabel ?? 'Enviar' }}</button>
<!--                              ^  null/undefined -> 'Enviar' -->
\end{lstlisting}

\section{\emoji{link} Property Binding \texttt{[ ]} y Event Binding \texttt{( )}}

\begin{lstlisting}
<!-- SIN corchetes — pasa el TEXTO literal "value()" -->
<input value="value()">  <!-- INCORRECTO -->

<!-- CON corchetes — evalua la EXPRESION -->
<input [value]="value()">  <!-- CORRECTO -->

<!-- Event binding — escucha eventos -->
<input (input)="onInput($event)">
<button (click)="onSubmit()">Enviar</button>
<input (keyup.enter)="buscar()">  <!-- solo al presionar Enter -->

<!-- Output de hijo -->
<app-select-field (valueChange)="state.setAnswer(field.id, $event)" />
\end{lstlisting}

\begin{tcolorbox}[regla, title={\emoji{pushpin} Regla}]
  \textbf{\texttt{[ ]}} → datos bajan del padre al hijo.\\
  \textbf{\texttt{( )}} → eventos suben del hijo al padre.
\end{tcolorbox}

\section{\emoji{performing-arts} \texttt{@if} / \texttt{@else} — Renderizado Condicional}

\begin{lstlisting}
<!-- VIEJO — *ngIf, el else era confuso -->
<div *ngIf="condicion; else miElse">Si</div>
<ng-template #miElse><div>No</div></ng-template>

<!-- NUEVO Angular 17+ -->
@if (condicion) {
  <div>Si</div>
} @else if (otraCondicion) {
  <div>Tal vez</div>
} @else {
  <div>No</div>
}
\end{lstlisting}

\subsection{En \texttt{survey-page.component.ts}}
\begin{lstlisting}
@if (loader.survey.isLoading()) {
  <mat-progress-bar mode="indeterminate" />
}

@if (loader.survey.error(); as err) {
  <!--                      ^ "as err" captura el valor en variable local -->
  <pre>{{ err }}</pre>
}

@if (loader.survey.value(); as schema) {
  <app-survey [schema]="schema" />
}
\end{lstlisting}

\begin{tcolorbox}[peligro, title={\emoji{warning} \texttt{@if} NO es \texttt{display:none}}]
\begin{center}
\begin{tabular}{p{5.5cm} p{6cm}}
\textbf{\texttt{display:none}} & \textbf{\texttt{@if}} \\
\hline
Elemento \textbf{existe} en el DOM & Elemento \textbf{no existe} \\
Angular ejecuta su código & Angular no ejecuta nada \\
Consume memoria y CPU & Cero costo cuando es \texttt{false} \\
\end{tabular}
\end{center}
\end{tcolorbox}

\section{\emoji{repeat-button} \texttt{@for} — Iteración Reactiva}

\begin{lstlisting}
<!-- VIEJO — trackBy era opcional -->
<div *ngFor="let s of sections; trackBy: trackByFn">{{ s.title }}</div>

<!-- NUEVO — track es OBLIGATORIO -->
@for (section of schema().sections; track section.id) {
  <section>
    <h3>{{ section.title }}</h3>
    @for (field of section.fields; track field.id) {
      <app-field-renderer [field]="field" ... />
    }
  </section>
}
\end{lstlisting}

\begin{tcolorbox}[analogia, title={\emoji{id-button} Por qué \texttt{track} importa}]
  Sin \texttt{track}: Angular destruye y recrea \textbf{todos} los elementos al cambiar la lista.\\
  Con \texttt{track}: Angular solo destruye \textbf{los que desaparecieron}.\\
  Animaciones intactas, focus conservado, inputs sin limpiarse.
\end{tcolorbox}

\subsection{Variables automáticas de \texttt{@for}}
\begin{lstlisting}
@for (item of lista; track item.id;
      let i = $index, first = $first, last = $last) {
  <!--      ^              ^               ^
            0,1,2...       true si primero   true si ultimo
            Otras: $even, $odd, $count -->
  <div [class.ultimo]="last">{{ i + 1 }}. {{ item.nombre }}</div>
}
\end{lstlisting}

\subsection{\texttt{@empty}}
\begin{lstlisting}
@for (encuesta of encuestas; track encuesta.id) {
  <tarjeta [data]="encuesta" />
} @empty {
  <p>No hay encuestas disponibles.</p>
}
\end{lstlisting}

\section{\emoji{left-right-arrow} \texttt{@switch} — Despacho por Tipo}

\begin{lstlisting}
@let f = field(); <!-- @let = variable local, evita llamar field() 7 veces -->

@switch (f.type) {
  @case ('text')     { <app-text-field     [field]="f" [value]="stringValue()"
                         (valueChange)="valueChange.emit($event)" /> }
  @case ('email')    { <app-text-field     [field]="f" [value]="stringValue()"
                         (valueChange)="valueChange.emit($event)" /> }
  @case ('number')   { <app-number-field   [field]="f" [value]="numberValue()"
                         (valueChange)="valueChange.emit($event)" /> }
  @case ('select')   { <app-select-field   [field]="f" [value]="stringValue()"
                         (valueChange)="valueChange.emit($event)" /> }
  @case ('radio')    { <app-radio-field    [field]="f" [value]="stringValue()"
                         (valueChange)="valueChange.emit($event)" /> }
  @case ('checkbox') { <app-checkbox-field [field]="f" [value]="arrayValue()"
                         (valueChange)="valueChange.emit($event)" /> }
}
\end{lstlisting}

\section{\emoji{hourglass-not-done} \texttt{@defer} — Lazy Loading en Template}

\begin{lstlisting}
@defer (on interaction) {
  <editor-gigante [contenido]="post" />   <!-- descarga al hacer click -->
} @placeholder {
  <div>Haz clic para editar...</div>      <!-- antes de interactuar -->
} @loading (minimum 300ms) {
  <mat-spinner diameter="32" />           <!-- mientras descarga -->
} @error {
  <p>Error al cargar el editor.</p>       <!-- si la descarga falla -->
}
\end{lstlisting}

\subsection{Triggers}
\begin{center}
\renewcommand{\arraystretch}{1.5}
\begin{tabular}{>{\ttfamily}p{4.5cm} p{8cm}}
\toprule
Trigger & Cuándo carga \\
\midrule
on idle & Cuando el browser no tiene nada que hacer \\
on viewport & Cuando el elemento entra en la pantalla \\
on interaction & Al hacer hover o click \\
on hover & Solo al pasar el mouse encima \\
on timer(3s) & Después de 3 segundos \\
when condicion() & Cuando una expresión Signal es \texttt{true} \\
\bottomrule
\end{tabular}
\end{center}

\section{\emoji{abacus} Pipes — Transformadores en el Template}

\begin{lstlisting}
{{ fecha   | date:'dd/MM/yyyy' }}  <!-- 2026-05-13 -> "13/05/2026" -->
{{ precio  | currency:'MXN'    }}  <!-- 1500 -> "MX$1,500.00" -->
{{ texto   | uppercase         }}  <!-- "hola" -> "HOLA" -->
{{ objeto  | json              }}  <!-- { id:1 } -> '{"id":1}' (debug) -->

<!-- Pipes encadenados izquierda a derecha -->
{{ nombre | uppercase | slice:0:10 }}  <!-- "Aldo Zetina" -> "ALDO ZETIN" -->
\end{lstlisting}

\section{\emoji{clipboard} Resumen — Bloque 2}

\begin{center}
\renewcommand{\arraystretch}{1.6}
\begin{tabular}{>{\ttfamily\bfseries}p{4.5cm} p{8cm}}
\toprule
Sintaxis & Para qué \\
\midrule
\{\{ expr \}\} & Mostrar texto (interpolación) \\
{[}prop{]} & Pasar dato a elemento o componente hijo \\
(evento) & Escuchar evento de elemento o componente hijo \\
{[}(modelo){]} & Two-way binding \\
@if / @else & Renderizar si condición es \texttt{true} \\
@for + track & Iterar lista de forma eficiente \\
@empty & Cuando la lista está vacía \\
@switch / @case & Seleccionar según valor \\
@let var = expr & Variable local en template \\
@defer (trigger) & Carga diferida de componentes pesados \\
pipe | nombre & Transformar valor sin modificar el original \\
\bottomrule
\end{tabular}
\end{center}

% ============================================================
\vfill
\begin{center}
  \textcolor{gray}{\small \emoji{graduation-cap} ApuntesAngular --- Mayo 2026}\\[0.2cm]
  \textcolor{gray}{\small Survey Maker --- Angular 19 Signal-First}\\[0.2cm]
  \textcolor{gray}{\small \texttt{github.com/AldoZM/Survey-Maker}}
\end{center}

\end{document}
